\documentclass[10pt,a4paper]{article}
\usepackage[margin=1.8cm,top=1.55cm,bottom=1.55cm]{geometry}
\usepackage{fontspec,xcolor,enumitem,titlesec,hyperref,tabularx,parskip}
\setmainfont{TeX Gyre Termes}
\definecolor{navy}{HTML}{183650}
\hypersetup{colorlinks=true,urlcolor=navy}
\setlength{\parindent}{0pt}
\setlength{\parskip}{3pt}
\setlist[itemize]{leftmargin=1.35em,topsep=2pt,itemsep=1pt,parsep=0pt}
\titleformat{\section}{\large\bfseries\color{navy}}{}{0pt}{}[\vspace{-4pt}\color{navy}\titlerule]
\titlespacing*{\section}{0pt}{10pt}{6pt}
\newcommand{\entry}[3]{\textbf{#1}\hfill{\small #2}\\{\emph{#3}}\par}
\begin{document}
\begin{center}
{\LARGE\bfseries Shu Yiyang}\\[4pt]
{\normalsize Undergraduate Student in Mathematics}\\[2pt]
School of the Gifted Young, University of Science and Technology of China\\[2pt]
Hefei, China \quad\textbar\quad \href{mailto:shuyiy@mail.ustc.edu.cn}{shuyiy@mail.ustc.edu.cn} \quad\textbar\quad \href{https://nicholasshu.com}{My Website}
\end{center}
\vspace{-7pt}
\section*{Research Interests}
Differential geometry, geometric analysis, convex geometry, and Kakeya-type problems. Particularly interested in convex Kakeya sets and related variants, extremal geometric configurations, and relationships among convexity, geometric constraints, area, diameter, and perimeter.
\section*{Education}
\entry{University of Science and Technology of China (USTC)}{Sep 2025--Present}{School of the Gifted Young; Undergraduate studies in mathematics (pure mathematics track)}
Admitted through the Special Class for Gifted Young Students selection process with Category A qualification.
\entry{Anqing No. 1 High School}{2023--2025}{Anqing, Anhui, China}
Second Prize in the Anhui provincial high school mathematics competition in both 2023 and 2024; ranked among the national top 20 in mathematics in the 2023 Xi'an Jiaotong University gifted-young-students selection examination.
\section*{Research Experience}
\entry{Research on Convex Kakeya Sets}{Sep 2026--Present}{Advisor: Prof. Zuoqin Wang, USTC}
\begin{itemize}
\item Investigating convex sets containing unit line segments in every direction and quantitative relationships among area, diameter, and perimeter.
\item Exploring extremal configurations, geometric inequalities, and limiting behavior of convex bodies.
\item Studying supporting hyperplanes, support functions, convex hulls, and related methods of convex geometry.
\end{itemize}
\section*{Independent Mathematical Manuscript}
\textbf{Some Geometric Properties of Hyperbolic Plane}\hfill{\small 2026}\\
\emph{Independent manuscript; Published in Issue 68 of \textit{Wa Ming}, a student-run publication of the School of Mathematical Sciences at USTC.}
\begin{itemize}
\item Studied the Poincaré upper half-plane model of hyperbolic geometry, including its compatibility with Hilbert's axiomatic system and its transformation groups.
\item Investigated properties of hyperbolic triangles, including the hyperbolic laws of sines and cosines, as well as Desargues' theorem and Pappus' theorem.
\end{itemize}
\section*{Honors and Awards}
\begin{itemize}
\item \textbf{2026:} First Place, USTC Robogame Competition.
\item \textbf{2025:} Provincial Second Prize, CMC Mathematics Competition.
\item \textbf{2025:} Category A qualification, USTC Special Class for Gifted Young Students admission.
\item \textbf{2023, 2024:} Second Prize (each year), Anhui provincial high school mathematics competition.
\item \textbf{2023:} National top 20 in mathematics, Xi'an Jiaotong University gifted-young-students selection examination.
\item \textbf{2023:} Successfully completed the Yau Mathematical Sciences Summer Camp.
\item \textbf{2019, 2020:} Second Prize (each year), provincial informatics competitions.
\item Additional high-school honors for student leadership and Communist Youth League service.
\item \textbf{2025:} Third Prize, USTC Jiuzhang Cup.
\end{itemize}
\section*{Mathematical Preparation}
\textbf{Foundations:} Mathematical analysis, linear algebra, ordinary differential equations, classical differential geometry, convex analysis.\\
\textbf{Advanced reading / study:} Smooth manifolds, differential topology, geometric analysis, partial differential equations, minimal surfaces, differential forms, de Rham cohomology.\\
\emph{Note: The list combines coursework and independent study; individual course grades are not specified.}
\section*{Leadership and Skills}
\textbf{Service:} Class Communist Youth League Branch Secretary (2025--Present); Student Union, Academic and Practice Department member (2026--Present).\\
\textbf{Technical:} LaTeX, XeLaTeX, Beamer, SolidWorks. \quad \textbf{Languages:} Mandarin Chinese (native); English (level to confirm).
\end{document}
